% Generated draft: requires mathematical and source review.
\documentclass[11pt]{article}
\usepackage[margin=25mm]{geometry}
\usepackage{fontspec}
\setmainfont{Latin Modern Roman}
\usepackage{amsmath,amssymb,mathtools,unicode-math}
\setmathfont{Latin Modern Math}
\usepackage{enumitem,needspace,xurl,hyperref}
\setlength{\parindent}{0pt}
\setlength{\parskip}{5pt}
\setlength{\emergencystretch}{3em}
\newcommand{\N}{\mathbb{N}}
\newcommand{\Z}{\mathbb{Z}}
\newcommand{\Q}{\mathbb{Q}}
\newcommand{\R}{\mathbb{R}}
\newcommand{\C}{\mathbb{C}}
\newcommand{\F}{\mathbb{F}}
\newcommand{\A}{\mathbb{A}}
\newcommand{\PP}{\mathbb{P}}
\newcommand{\E}{\mathbb{E}}
\newcommand{\eps}{\varepsilon}
\newcommand{\dd}{\,\mathrm{d}}
\begin{document}
\section*{UnsolvedMath ID: 9700026}
\textbf{Source problem code:} catalogue entry
\subsection{Definitions and mathematical statement}
Consider the source's associated influence-cell dynamical system with fixed geometric data in general position and strictly positive initial weights. Write its weights as $w_1(t),\ldots,w_N(t)$ and retain its parameters $\alpha,\beta$; the actual evolution and normalization are not supplied.

The first claimed regime is
\[
\alpha>1\quad\text{or}\quad0<\beta<2\alpha.
\]
In this regime the conjecture requires existence of an index $i$ such that $w_i(t)\to1$ as time tends to infinity. The winning index may depend on the initial values; no fixed universal winner is asserted. If the source normalizes the weights to sum to one, the remaining weights then tend to zero, but that normalization must be recovered from its system.

In the second regime,
\[
0<\alpha<1,\qquad \beta>2\alpha,
\]
all weights should converge to positive limits. Independence of initial values means that, with the geometric data and parameters fixed, there is one positive limiting vector $w^*$ such that every permitted positive initial vector yields $w_i(t)\to w_i^*$ for every $i$. It does not mean independence from the city locations or parameters, and it is not an almost-every-initial-value claim.

No conclusion is stated for the boundary cases $\alpha=1$ or $\beta=2\alpha$. The record also does not give the full parameter domain outside the displayed inequalities, the meaning of general position, the influence cells, or the differential/iterated equations. Thus these quantifiers preserve the exact dichotomy while the missing model prevents a self-contained stability theorem.
\subsection{Short English statement}
Prove the source city dynamics either concentrates at one weight in the stated unbalanced regimes or converges to a positive initial-condition-independent vector in the stated balanced regime. The evolution equations and boundary cases are not supplied.
\subsection{Sources}
\begin{itemize}
\item Local UnsolvedMath registry, permanent ID 9700026, exact source code catalogue entry, ``Stability dichotomy for the associated city dynamical system''; attributed to David Aldous, no proposed year supplied. Registry status open, classification OPEN-TRIAGE.
\item David Aldous, \emph{Open problems}, item \texttt{cities.html}, technical notes, Conjecture 32, \url{https://www.stat.berkeley.edu/~aldous/Research/OP/index.html}. Supplied HTML-text with linked-primary-PDF-text extraction accessed July 29, 2026. Provenance records 23 active pages with compound bullets retained; no linked PDF was accessed here.
\item The supplied summary names the Aldous--Huang city model and \texttt{arXiv:1209.5120}, without a title or exact equation reference, and reports no proof of the dichotomy or verified follow-up resolution. It does not supply the missing dynamics.
\item Provenance flags current primary-literature status and redistribution terms NEEDS\_REVIEW.
\end{itemize}
\end{document}
